
\documentclass[10pt,letterpaper]{article}

\usepackage[margin=0.82in]{geometry}
\usepackage{times}
\usepackage{microtype}
\usepackage{enumitem}
\usepackage{amsmath}
\usepackage{booktabs}

\setlength{\parindent}{0pt}
\setlength{\parskip}{4pt}
\setlist{nosep,leftmargin=1.4em}
\setcounter{secnumdepth}{2}

\title{\vspace{-1.5em}\textbf{SPECTRE: An Event-Driven Vision-Language\\
Tactical Coach for Apex Legends}}
\author{AI Systems Architecture \\ Utah State University}
\date{}

\begin{document}
\maketitle
\vspace{-2em}

\section{Introduction and Motivation}

Competitive multiplayer games provide a challenging environment for
intelligent decision support. In Apex Legends, players must continuously
interpret their surroundings, track limited resources, coordinate with
teammates, and make tactical decisions under strict time constraints.
The outcome of an engagement often depends less on mechanical execution
alone than on the timing and context of individual decisions. A player
may lose an otherwise favorable fight by abandoning cover, failing to
recognize an isolated opponent, or engaging without sufficient resources.

Although gameplay recordings enable retrospective analysis, extracting
meaningful lessons from them remains difficult. Conventional performance
statistics summarize outcomes through kills, damage, and placement, but
provide limited insight into the decisions that produced those outcomes.
A useful coaching system must instead reconstruct the circumstances of
an engagement, identify the information available to the player, and
evaluate whether alternative decisions could have improved the result.

Vision-language models (VLMs) offer a promising foundation for this
task because they can jointly interpret visual scenes and generate
natural-language explanations. However, competitive gameplay presents
several difficulties for direct VLM analysis. Important information is
distributed across small HUD elements, game mechanics are highly
specialized, and meaningful events may unfold between consecutive
sampled frames. Treating every frame as an independent image can
therefore produce incomplete or misleading tactical interpretations.

We propose \textbf{SPECTRE}, an offline tactical coaching system that
combines conventional computer vision with temporally organized
vision-language reasoning. Its central premise is that reliable
extraction of structured gameplay information can reduce the perceptual
burden on a VLM, allowing the model to concentrate on spatial context,
decision-making, and event interpretation.

\section{Problem Statement and Objectives}

The problem addressed by SPECTRE is the conversion of raw gameplay
recordings into evidence-based tactical feedback. This requires more
than object recognition: the system must preserve relevant observations
over time and distinguish visible facts from uncertain interpretations.

The project has four objectives:

\begin{enumerate}
    \item \textbf{Structured perception:} Extract relevant HUD information,
    including squad composition, defensive resources, equipped weapons,
    and interface states.
    \item \textbf{Temporal organization:} Associate observations with
    timestamps and maintain a consistent representation of changing
    gameplay conditions.
    \item \textbf{Contextual interpretation:} Combine gameplay imagery
    with structured observations to identify engagements, positioning,
    and tactical decisions.
    \item \textbf{Actionable feedback:} Produce concise, timestamped
    assessments that explain observed decisions and identify defensible
    alternatives.
\end{enumerate}

The initial scope is offline match analysis. Real-time coaching is
considered a longer-term extension rather than a requirement of the
first prototype.

\section{Proposed Architecture}

SPECTRE separates the system into three principal stages.

\subsection{Visual Perception}

Recorded gameplay is sampled at configurable intervals. Specialized
computer-vision operations process predefined interface regions using
template matching, color analysis, and optical character recognition
where appropriate. Screen-state identification determines whether the
player is in normal gameplay, inventory, map, spectating, or another
interface state, preventing unrelated interface elements from being
interpreted as gameplay measurements.

This division is motivated by efficiency and reliability. Stable HUD
elements are comparatively well suited to deterministic processing,
whereas spatial relationships and tactical context require more
flexible visual reasoning.

\subsection{Temporal Representation}

The extracted information is organized into a chronological match
record. Each observation contains a timestamp, corresponding video
frame, recognized screen state, and available HUD measurements.
Information that is not visible is distinguished from information
known to be absent. This distinction is important when the player
changes weapons, opens an inventory, or temporarily obscures the HUD.

The resulting representation provides the basis for identifying
changes in player resources and reconstructing gameplay events across
multiple observations.

\subsection{Vision-Language Tactical Analysis}

A VLM receives ordered gameplay frames together with the corresponding
structured observations. The model is instructed to identify visible
events, assess player positioning and resource management, and generate
recommendations supported by the supplied evidence.

The analysis explicitly distinguishes observed events from uncertain
interpretations. For example, a visible knockdown must not be treated
as confirmation that an entire opposing squad has been eliminated.

\section{Implementation and Evaluation Plan}

The prototype will be implemented in Python using OpenCV for video
processing and HUD extraction, JSON for structured observations, and a
locally hosted VLM for multimodal interpretation. The architecture
will preserve intermediate outputs so that perception and reasoning
can be evaluated independently.

Evaluation will address three questions:

\begin{itemize}
    \item \textbf{Perception:} How reliably are HUD elements identified
    across different gameplay and interface states?
    \item \textbf{Temporal understanding:} Can the system reconstruct
    meaningful events from sampled frames without confusing missing
    observations with actual changes?
    \item \textbf{Coaching quality:} Does structured HUD context help
    produce more specific and evidence-supported recommendations than
    analysis based on gameplay imagery alone?
\end{itemize}

Representative match segments will be inspected manually to identify
correct observations, missed events, and unsupported conclusions.
The initial evaluation emphasizes functional integration and
qualitative failure analysis; larger-scale quantitative benchmarking
will follow if time permits.

\section{Expected Outcomes}

The expected outcome is a working offline prototype that accepts
recorded Apex Legends gameplay and produces a structured tactical
report with timestamped events, observations, and recommendations.

The broader contribution is architectural: SPECTRE investigates how
deterministic perception, temporal context, and multimodal reasoning
can be combined to support intelligent analysis of a dynamic
environment. Future work may incorporate game-specific knowledge,
fight-aware sampling, supervised fine-tuning, and event-driven
inference under real-time deadlines.

\end{document}
